% Centri — Weekly progress (week ending 2026-07-21)
% Standalone deck: every slide readable cold, no narration needed.
% Compile: pdflatex -interaction=nonstopmode centri-weekly-2026-07-21.tex  (run twice)
\documentclass[aspectratio=169,10pt]{beamer}
\usetheme{metropolis}
\metroset{sectionpage=progressbar, subsectionpage=none, numbering=fraction, progressbar=frametitle}

\usepackage{booktabs}
\usepackage{amsmath,amssymb}
\usepackage{xcolor}
\usepackage{colortbl}
\usepackage{graphicx}
\graphicspath{{figs/}}
\usepackage[most]{tcolorbox}

% ---- palette (shared with the other Centri decks) ------------------------
\definecolor{cMeas}{HTML}{1F6FB2}
\definecolor{cPed}{HTML}{B5651D}
\definecolor{cApp}{HTML}{2E7D32}
\definecolor{cInfra}{HTML}{5E35B1}
\definecolor{cAccent}{HTML}{C62828}
\definecolor{cGrey}{HTML}{616161}
\setbeamercolor{frametitle}{bg=cMeas}
\setbeamercolor{progress bar}{fg=cAccent}

\newcommand{\ok}{{\color{cApp}\textbf{\checkmark}}}
\newcommand{\nope}{{\color{cGrey}--}}
\newcommand{\todo}{{\scriptsize\color{cAccent}\textbf{[to fill in]}}}
% one-line, self-contained takeaway centred under the slide body
\newcommand{\takeaway}[1]{\vfill{\footnotesize\itshape\color{cGrey}#1\par}}

\title{Centri --- Measurements That Can Be Trusted}
\subtitle{Video-generated tiered physics material, with every number checked against the video it came from}
\author{Damar Albaribin}
\institute{MSc Thesis --- advisor: Prof.\ Wu-Yuin Hwang}
\date{Week ending 2026-07-21}

\begin{document}
\maketitle

% =========================================================================
\begin{frame}{What Centri produces, and where the risk sits}
\small
\textbf{Centri} takes an ordinary phone video of an object spinning in a circle and, using
computer-vision tracking (\emph{no} phone sensors), turns it into physics learning material at
\textbf{three levels} --- basic, intermediate, advanced --- with \textbf{annotated figures}: the
tracked frame, the motion graphs, a data table.
\vspace{0.4em}

\textbf{P-MAGIC} is the sibling app: same three-level multimodal material, but measured from a
phone's motion \emph{sensors}. Ten physics teachers rated its material \textbf{8.2--8.6\,/\,10}.
A sensor reports the motion directly. \textbf{A video does not} --- every quantity must first be
recovered from pixels, and that recovery can be wrong in ways the finished material cannot show.
\vspace{0.4em}

\begin{tcolorbox}[colback=cMeas!6,colframe=cMeas,boxsep=2pt,left=6pt,right=6pt,top=3pt,bottom=3pt]
\small So the material is only as good as the measurement underneath it. Centri now
\textbf{corrects for the angle the video was shot from}, \textbf{checks each measurement against
the video it came from}, and \textbf{shows the measurement in the figures} rather than an idealised
drawing of it.
\end{tcolorbox}
\takeaway{Each of those three claims is unpacked on its own slide, with the measurement that supports it.}
\end{frame}

% =========================================================================
\begin{frame}{The camera angle is corrected before anything is measured}
\scriptsize
A circle filmed from off to one side photographs as an \textbf{oval}, with its near side looking
bigger than its far side, so an object going round at a \emph{steady} rate appears to speed up and
slow down once per lap. No smoothing can remove that: it repeats identically every lap, exactly as a
real rhythm would.
\vspace{1pt}
\begin{center}
\includegraphics[width=0.985\textwidth]{rect4046_slide.png}
\end{center}
\vspace{1pt}
One fact of perspective undoes it: a circle's \textbf{true centre} and the \textbf{centre of the oval
you see} coincide only when the camera looks straight on, so their separation measures the slant ---
with no camera calibration and nothing known about the room.
\takeaway{It changes where \emph{within} a lap the object is judged to be, not how many laps happened
--- the average turn rate moves 0.2\%.}
\end{frame}

% =========================================================================
\begin{frame}{The same video, measured two ways}
\scriptsize
Both pictures are real output from the same 15-second clip of a playground wheel. The difference is
entirely in how the measurement was made --- naively on the left, corrected on the right.
\vspace{2pt}
\begin{columns}[T]
\begin{column}{0.30\textwidth}
\centering
\includegraphics[width=0.80\textwidth]{wheel_naive_crop.png}\\[1pt]
{\scriptsize\textbf{\color{cAccent}Measured naively.} The blue dot, meant to mark the handle,
floats clear of it; the ring is a circle on a path that is not one.}
\end{column}
\begin{column}{0.30\textwidth}
\centering
\includegraphics[width=0.80\textwidth]{wheel_corrected_crop.png}\\[1pt]
{\scriptsize\textbf{\color{cApp}Measured correctly.} The dot sits on the handle,
and the ring follows the oval the wheel really traces.}
\end{column}
\begin{column}{0.38\textwidth}
\vspace{6pt}
\renewcommand{\arraystretch}{1.26}
\begin{tabular}{@{}lrr@{}}
\toprule
& \color{cAccent}naive & \color{cApp}correct \\
\midrule
radius (m)              & 0.281 & \textbf{0.260} \\
turn rate (rad/s)       & 7.88 & \textbf{7.80} \\
inward accel.\ (m/s\textsuperscript{2}) & 22.8 & \textbf{16.0} \\
largest accel.          & 235.5 & \textbf{23.0} \\
radius wobble (m)       & 0.114 & \textbf{0.009} \\
one lap (s)             & 1.07 & \textbf{0.81} \\
\bottomrule
\end{tabular}
\end{column}
\end{columns}
\takeaway{The left column is not noisy, it is \emph{confidently wrong} --- its lap time contradicts its
own turn rate (7.88\,rad/s means 0.80\,s, not 1.07\,s). No finished worksheet would reveal that, which
is why the checks run on the measurement rather than on the prose.}
\end{frame}

% =========================================================================
\begin{frame}{Each measurement is checked against the video it came from}
\small
Four automatic checks, each rejecting one specific way a plausible-looking number can be wrong.
They cost nothing, need no AI model, and run before any material is written.
\vspace{0.3em}
\scriptsize
\renewcommand{\arraystretch}{1.5}
\begin{tabular}{@{}p{5.0cm}p{7.6cm}@{}}
\toprule
\textbf{The check} & \textbf{The mistake it refuses to ship} \\
\midrule
the object's path and the centre it turns about are measured on the \textbf{same picture} &
the two arriving from differently cropped copies of the frame --- which reads as a genuinely
wobbly wheel, so the warnings blame the physics instead of the arithmetic \\
the slant correction is applied \textbf{only where a slant was measured} &
``correcting'' a video already shot straight on, which would fit the correction to noise and
inject an error rather than remove one \\
the figures plot \textbf{measured positions} &
a marker drawn where the idealised circle says the object should be, rather than where it
actually was \\
worked examples \textbf{close on their own printed numbers} &
a worked example whose stated answer is not what its own printed multiplication gives \\
\bottomrule
\end{tabular}
\takeaway{Each check earns its place by having caught a real error on real footage --- the first two
found faults that had already survived a full review of the finished material.}
\end{frame}

% =========================================================================
\begin{frame}{The evidence set: seven clips, four phenomena}
\scriptsize
Every clip Centri is willing to teach from, and what is actually \emph{measured} in each. A clip
earns a place here only when its numbers can be defended, not when the footage looks good.
\vspace{3pt}
\renewcommand{\arraystretch}{1.35}
\begin{center}
\begin{tabular}{@{}lcccc@{}}
\toprule
& \textbf{position of the} & \textbf{real-world} & \textbf{slant} & \textbf{material passed} \\
& \textbf{object, every frame} & \textbf{size} & \textbf{correction} & \textbf{every check} \\
\midrule
playground wheel & \ok & \ok & \ok\ {\color{cGrey}applied} & \ok \\
turntable {\color{cGrey}\,(3 clips)} & \ok & \ok & \nope\ {\color{cGrey}shot straight on} & \ok \\
computer fan & \ok & \ok & \nope\ {\color{cGrey}shot straight on} & \ok \\
ceiling fan {\color{cGrey}\,(2 clips)} & \nope\ {\color{cGrey}rate only} & \ok\ {\color{cGrey}tape-measured} & \nope\ {\color{cGrey}shot straight on} & \ok \\
\bottomrule
\end{tabular}
\end{center}
\vspace{4pt}
{\scriptsize The ceiling fan is the one exception, and it is labelled as such wherever it is used:
a spinning blade has no fixed point to follow, so its \emph{rate} and \emph{size} are measured but
its path is reconstructed from the rate. Its speed and acceleration are sound; its drawn circle is
a model, not a photograph of the motion.}
\takeaway{Exactly one clip needs the slant correction --- the one that was actually filmed at a slant.
The others are left untouched, and record why.}
\end{frame}

% =========================================================================
\begin{frame}{What the learner actually receives}
\scriptsize
Real output from the same run --- the figures that go into the student worksheet alongside the prose.
\vspace{3pt}
\begin{columns}[T]
\begin{column}{0.52\textwidth}
\centering
\includegraphics[width=\textwidth]{wheel_graph.png}\\[2pt]
{\scriptsize\textbf{Turn rate over time.} The shaded band names the motion phase the graph actually
draws --- here a single long \emph{slowing down} --- and the dashed line is the clip average, so a
frozen average is never mistaken for an instant reading.}
\end{column}
\begin{column}{0.44\textwidth}
\vspace{4pt}
{\scriptsize The remaining ripple in this line is \emph{real}: the handle is a solid knob roughly
110 pixels across, and its centre in the picture genuinely shifts as different faces of it come into
view. That is the subject of the next slide.}
\vspace{6pt}

{\scriptsize Every value that reaches the worksheet is one of these measured quantities, tagged
\emph{(avg)} whenever it is a whole-clip average:}
\vspace{3pt}

\renewcommand{\arraystretch}{1.35}
\begin{tabular}{@{}lrl@{}}
\toprule
\textbf{measurement} & \textbf{value} & \\
\midrule
radius & 0.260 & m \\
turn rate {\color{cGrey}(avg)} & 7.80 & rad/s \\
inward acceleration {\color{cGrey}(avg)} & 16.05 & m/s\textsuperscript{2} \\
speed along the path & 2.03 & m/s \\
time for one lap & 0.81 & s \\
laps per second & 1.23 & Hz \\
\bottomrule
\end{tabular}
\end{column}
\end{columns}
\takeaway{Three levels of worksheet are written from these same six numbers, so the basic and advanced
versions can differ in demand without ever disagreeing on fact.}
\end{frame}

% =========================================================================
\begin{frame}{What limits accuracy now is the sticker, not the geometry}
\small
With the camera angle corrected, the remaining scatter was measured to find out what is still in
the way. The answer is not the geometry.
\vspace{0.4em}
\scriptsize
\renewcommand{\arraystretch}{1.4}
\begin{tabular}{@{}p{6.4cm}rp{5.6cm}@{}}
\toprule
\textbf{What was measured} & & \textbf{What it means} \\
\midrule
how much of the leftover error repeats each lap &
\textbf{1.3\%} & a wrong shape model would show up here. Nothing left to model. \\
how much the tracked handle's \emph{apparent size} changes as it goes round &
\textbf{75--82\%} & we see different faces of a solid knob at different distances \\
leftover scatter, against the size of the handle itself &
\textbf{$\approx$15\%} & ``where exactly is the handle?'' has no better answer than this \\
\bottomrule
\end{tabular}
\vspace{0.5em}

\small The object being followed is a \emph{lump}, not a point, so its centre in the picture shifts
as it turns. A better shape model cannot fix that --- but a \textbf{small, flat, high-contrast
sticker} can, and that is a change to how clips are filmed, not to the software. The same fix
already rescued two earlier clips.
\takeaway{Worth knowing before building anything more elaborate: the next real gain comes from the
camera work, not from more mathematics.}
\end{frame}

% =========================================================================
\begin{frame}{A number that could not be true, and where it came from}
\small
One clip shows a turntable \textbf{speeding up} while it is slowing down --- nothing pushes it, so
that cannot happen. Rather than smooth it away, it was traced back stage by stage.
\vspace{0.3em}
\scriptsize
\renewcommand{\arraystretch}{1.3}
\begin{tabular}{@{}p{4.3cm}p{4.0cm}p{5.2cm}@{}}
\toprule
\rowcolor{cMeas!12}
\textbf{Stage checked} & \textbf{What was found} & \textbf{Why it matters} \\
\midrule
The camera, while filming &
\textbf{one frame missing}, at 2.99\,s &
the phone skipped a beat; its own clock records the gap \\
Turning video into still pictures &
the gap is filled by \textbf{repeating the frame before it} &
a valid image, so the object appears not to have moved \\
Measuring speed from positions &
timing assumed \textbf{perfectly even} &
one motionless step, then a normal one, reads as ``stopped, then sped up'' \\
\bottomrule
\end{tabular}
\vspace{0.5em}

\scriptsize \textbf{\color{cApp}Now fixed:} the software reads the video's own clock and discards
positions taken from a repeated picture. It cannot be done the obvious way --- spotting an object
that stopped moving --- because every clip ends with the object genuinely at rest.
\vspace{0.2em}

\scriptsize\color{cGrey}Scope, stated plainly: this explains about \textbf{one-eighth} of that
clip's anomaly, and one frame in 12{,}900 is affected. The rest stays unexplained, and the clip
carries that label rather than being presented as clean.
\takeaway{The discipline matters more than the bug: a number that violates physics is traced to a
cause, and whatever remains unexplained keeps its warning label.}
\end{frame}

% =========================================================================
\begin{frame}{Planned: filming the ceiling fan with sensors running}
\small
The open question this project has not yet answered directly: \textbf{does a video measurement agree
with a sensor measurement of the same motion?} Every comparison so far is against a different setup
on a different day.
\vspace{0.3em}

\begin{tcolorbox}[colback=cPed!6,colframe=cPed,boxsep=1pt,left=6pt,right=6pt,top=2pt,bottom=2pt]
\small \textbf{The plan.} Film the ceiling fan while a phone strapped to the blade records its own
motion sensors. One motion, two independent instruments, same moment --- so the two turn-rate curves
can be laid on top of each other.
\end{tcolorbox}
\vspace{0.3em}

\scriptsize
\renewcommand{\arraystretch}{1.35}
\begin{tabular}{@{}p{5.3cm}p{5.6cm}c@{}}
\toprule
\textbf{To settle} & \textbf{Why it matters} & \textbf{result} \\
\midrule
how closely the two agree on turn rate & the headline number for the video-vs-sensor claim & \todo \\
whether the video keeps up during speed-up & the hardest part of the motion to catch & \todo \\
what the phone's own mass does to the fan & the instrument changes what it measures & \todo \\
which one is better on real-world size & the video needs a ruler; the sensor does not & \todo \\
\bottomrule
\end{tabular}
\takeaway{This is the measurement that would let Centri claim video is a valid substitute for sensors,
rather than merely a convenient one.}
\end{frame}

% =========================================================================
\begin{frame}{Also planned, and open decisions}
\small
\begin{columns}[T]
\begin{column}{0.5\textwidth}
\textbf{\color{cMeas}Measurement}
\begin{itemize}\setlength\itemsep{2pt}
  \item \textbf{Re-film the park exercise wheel} --- the best motion in the collection (27 turns,
        the only one that coasts fully to rest) but nothing can currently follow it. Needs one strip
        of bright tape on a handle. \todo
  \item \textbf{A ruler in every scene.} Real-world size is the one quantity a video cannot get on
        its own, and it is the weakest link in two clips. \todo
  \item \textbf{One unexplained clip} (bicycle wheel) --- a repeating wobble that is \emph{not} the
        camera angle. Cause unknown. \todo
\end{itemize}
\end{column}
\begin{column}{0.5\textwidth}
\textbf{\color{cPed}Teaching}
\begin{itemize}\setlength\itemsep{2pt}
  \item \textbf{Ask before showing?} A ``what do you predict?'' box before the measurement is
        revealed. \todo
  \item \textbf{A compare-two task?} P-MAGIC has students compare two experiments; from one video
        we cannot --- pair two clips, or compare two phases of one? \todo
  \item \textbf{Teacher ratings.} Give the rating sheets to real teachers and check how well their
        scores agree with the automatic ones --- the step that turns these results from preliminary
        into publishable. \todo
\end{itemize}
\end{column}
\end{columns}
\takeaway{Placeholder slide --- each item above needs a decision or a number before it can be
presented as a result.}
\end{frame}

\end{document}
